\chapter{Ivanti: intrusión mediante una cadena de fallos}
\label{chap:incidente2}

\section{Contexto y justificación de selección}
El 5 de octubre de 2026, Volexity inició una investigación
que identificó actividad maliciosa desde el día 3 en una VPN
Ivanti Connect Secure 22.8R2.4, actualizada en aquel momento.
El atacante, denominado UTA0576, encadenó vulnerabilidades
para ejecutar comandos como administrador sin autenticarse.
Se selecciona porque muestra cómo comprometer una herramienta
de acceso remoto puede abrir el camino hacia los sistemas
que debería proteger~\cite{volexity2026ivanti}.

\section{Desarrollo del incidente y CVE/CVSS}
El atacante eludió los controles de autenticación y alcanzó
funciones internas de carga de archivos. Aprovechó el intervalo
entre la escritura y eliminación de un archivo rechazado para
ejecutar código. Después manipuló datos de configuración
utilizados por un servicio privilegiado y obtuvo permisos
\textit{root}~\cite{volexity2026ivanti}.

Los registros oficiales permiten distinguir tres componentes:
\textbf{CVE-2026-104935}, evasión de autenticación que permite
ejecutar comandos remotos, con \textbf{CVSS v3.1: 8,6};
\textbf{CVE-2026-118527}, nombres predecibles que permiten
depositar archivos en ubicaciones conocidas, con
\textbf{CVSS v3.1: 4,3}; y \textbf{CVE-2026-111357},
escalada de privilegios mediante confianza indebida en la
configuración DSCache, valorada con \textbf{CVSS v3.1: 7,0}
en el registro de MITRE. Son puntuaciones publicadas para
fallos individuales, no una valoración conjunta de la
cadena~\cite{cve2026ivanti104935,cve2026ivanti118527,
    cve2026ivanti111357}.

El primer fallo permite acceso por red, sin privilegios
previos ni interacción, según su vector oficial:\par
\noindent{\small\ttfamily
    CVSS:3.1/AV:N/AC:L/PR:N/UI:N/S:U/C:H/I:L/A:L}\par
\noindent La combinación elimina las barreras que limitarían
la explotación aislada de los otros
componentes~\cite{cve2026ivanti104935}.

\section{Impacto y mitigación}
Volexity documentó extracción de credenciales, exfiltración
de metadatos de Active Directory y búsquedas de contraseñas
administrativas y claves de recuperación BitLocker.
El atacante también borró registros y modificó componentes
en memoria; el comprobador de integridad de Ivanti no
detectaba esas modificaciones~\cite{volexity2026ivanti}.
Para las organizaciones afectadas, las credenciales robadas
pueden facilitar accesos posteriores y suplantaciones; para
sus usuarios, implican riesgos sobre información y servicios.
Las fuentes no ofrecen un recuento global de personas
afectadas.

La respuesta requiere instalar una versión que corrija
todos los componentes, según la rama: el aviso canadiense
actualizado recoge 25.1.1.3, 25.1.3.1 y 22.8R2.6.
CISA incorporó el primer CVE a su catálogo el 13 de octubre
y los otros dos el día 30~\cite{canada2026ivanti}.
Además de actualizar, deben investigarse los accesos
desde la VPN y revocarse las credenciales comprometidas;
una comprobación de integridad limpia no descarta la
intrusión~\cite{volexity2026ivanti}.

\section{Análisis crítico}
La principal lección es que estar actualizado no elimina
el riesgo de vulnerabilidades desconocidas. Tampoco basta
con valorar cada fallo por separado: una debilidad moderada
puede resultar decisiva dentro de una cadena.
La defensa debe combinar actualización, limitación de
privilegios y vigilancia de los accesos internos.
La ausencia de alteraciones detectables en disco debe
interpretarse junto con otras evidencias, no como garantía
de seguridad.